\documentclass[a4paper]{article}

\usepackage[a4paper, left=3cm, right=2cm, top=2cm, bottom=2cm]{geometry}

% Phông chữ Times New Roman
\usepackage{fontspec}
\setmainfont{Times New Roman}

% Cách dòng đơn + cách đoạn 3pt
\usepackage{setspace}
\singlespacing
\setlength{\parskip}{3pt}
\setlength{\parindent}{1cm}

% Cỡ chữ 13
\renewcommand{\normalsize}{\fontsize{13pt}{15.6pt}\selectfont}
\normalsize

\usepackage{cite}
\usepackage{amsmath}
\usepackage{graphicx}
\usepackage{booktabs}
\usepackage{array}
\usepackage{url}
\usepackage{listings}
\usepackage{indentfirst}
\usepackage{float}
\usepackage{tikz}
\usetikzlibrary{
	arrows.meta,
	calc,
	shapes,
	shapes.misc,
	shapes.arrows,
	chains,
	matrix,
	positioning,
	scopes,
	decorations.pathmorphing,
	shadows
}

\graphicspath{{./}{paper/}{paper/figures/}}

% =====================================================================
% Định dạng theo mẫu báo cáo cuộc thi KHKT
% (trang bìa, MỤC LỤC, DANH MỤC HÌNH VẼ, DANH MỤC BẢNG BIỂU)
% =====================================================================
\usepackage{titlesec}
\usepackage{tocloft}
\usepackage[labelsep=period]{caption}

% Đánh số: PHẦN I, PHẦN II, ... và mục con I.1, I.2, ...
\renewcommand{\thesection}{\Roman{section}}
\renewcommand{\thesubsection}{\thesection.\arabic{subsection}}
\titleformat{\section}{\normalfont\Large\bfseries}{PHẦN \thesection:}{0.5em}{}
\titleformat{\subsection}{\normalfont\normalsize\bfseries}{\thesubsection.}{0.5em}{}

% Tên chú thích hình, bảng
\renewcommand{\figurename}{Hình}
\renewcommand{\tablename}{Bảng}

% Tiêu đề các danh mục
\renewcommand{\contentsname}{MỤC LỤC}
\renewcommand{\listfigurename}{DANH MỤC HÌNH VẼ, SƠ ĐỒ}
\renewcommand{\listtablename}{DANH MỤC BẢNG BIỂU}
\renewcommand{\refname}{TÀI LIỆU THAM KHẢO}

% Tiêu đề danh mục căn giữa, in đậm
\renewcommand{\cfttoctitlefont}{\hfill\normalsize\bfseries}
\renewcommand{\cftaftertoctitle}{\hfill\null}
\renewcommand{\cftloftitlefont}{\hfill\normalsize\bfseries}
\renewcommand{\cftafterloftitle}{\hfill\null}
\renewcommand{\cftlottitlefont}{\hfill\normalsize\bfseries}
\renewcommand{\cftafterlottitle}{\hfill\null}
\setlength{\cftbeforetoctitleskip}{0pt}
\setlength{\cftaftertoctitleskip}{6pt}
\setlength{\cftbeforeloftitleskip}{18pt}
\setlength{\cftafterloftitleskip}{6pt}
\setlength{\cftbeforelottitleskip}{18pt}
\setlength{\cftafterlottitleskip}{6pt}

% Mục lục: "PHẦN I: ...", in đậm, dấu chấm dẫn
\renewcommand{\cftsecfont}{\bfseries}
\renewcommand{\cftsecpagefont}{\bfseries}
\renewcommand{\cftsecpresnum}{PHẦN~}
\renewcommand{\cftsecaftersnum}{:}
\renewcommand{\cftsecleader}{\bfseries\cftdotfill{\cftdotsep}}
\setlength{\cftsecnumwidth}{6.5em}
\setlength{\cftbeforesecskip}{4pt}

\renewcommand{\cftsubsecfont}{\bfseries}
\renewcommand{\cftsubsecpagefont}{\normalfont}
\renewcommand{\cftsubsecaftersnum}{.}
\setlength{\cftsubsecindent}{1em}
\setlength{\cftsubsecnumwidth}{5.5em}

% Danh mục hình: "Hình 1. ..."
\renewcommand{\cftfigpresnum}{Hình~}
\renewcommand{\cftfigaftersnum}{.}
\setlength{\cftfigindent}{0pt}
\setlength{\cftfignumwidth}{4em}

% Danh mục bảng: "Bảng 1. ..."
\renewcommand{\cfttabpresnum}{Bảng~}
\renewcommand{\cfttabaftersnum}{.}
\setlength{\cfttabindent}{0pt}
\setlength{\cfttabnumwidth}{4em}

\begin{document}
	
	\title{Báo cáo tóm tắt kết quả nghiên cứu \\
		{\large MemoLang: Ngữ nghĩa đọc phá hủy cho các khối lượng công việc ma trận trên GPU}}
	
	\author{
		Independent Interpreter Project\\
		Tháng 9 năm 2026, Hà Nội}
	
	% ================= TRANG 1: TRANG BÌA =================
	\thispagestyle{plain}
	\begin{center}
		\bfseries
		CUỘC THI KHOA HỌC, KỸ THUẬT DÀNH CHO HỌC SINH TRUNG HỌC CẤP TRƯỜNG NĂM HỌC 2026 -- 2027
	\end{center}
	
	\vspace{5cm}
	
	\begin{center}
		\bfseries
		BẢN BÁO CÁO TÓM TẮT KẾT QUẢ NGHIÊN CỨU
	\end{center}
	
	\vspace{3cm}
	
	\begin{center}
		{\Large\bfseries\itshape Tên đề tài:}\\[10pt]
		{\Large\bfseries
		\MakeUppercase{MemoLang: Nghiên cứu và chế tạo ứng dụng cơ chế giải phóng dữ liệu đệm trong phát triển AI}\par}
	\end{center}
	
	\vspace{3cm}
	
	\noindent\textbf{\textit{Lĩnh vực:} Phần mềm hệ thống}
	
	\vspace{6pt}
	
	\vfill
	
	\begin{center}
		\textbf{\textit{Hà Nội, tháng 9 năm 2026}}
	\end{center}
	
	\newpage
	
	% ================= TRANG 2–3: MỤC LỤC + DANH MỤC =================
	\tableofcontents
	
	\listoffigures
	\listoftables
	
	% ----- DANH MỤC TỪ VIẾT TẮT -----
    \vspace{10cm}
	\vspace{18pt}
	\begin{center}
		\textbf{DANH MỤC TỪ VIẾT TẮT}
	\end{center}
	\vspace{-6pt}
	\begin{center}
		\renewcommand{\arraystretch}{1.15}
		\begin{tabular}{>{\bfseries}l p{6.2cm} p{6.2cm}}
			\toprule
			\textbf{Viết tắt} & \textbf{Tên đầy đủ} & \textbf{Ý nghĩa} \\
			\midrule
			AI      & Artificial Intelligence & Trí tuệ nhân tạo \\
			AOT     & Ahead-Of-Time (compilation) & Biên dịch trước khi chạy \\
			AST     & Abstract Syntax Tree & Cây cú pháp trừu tượng \\
			cuBLAS  & CUDA Basic Linear Algebra Subprograms & Thư viện đại số tuyến tính cơ bản trên CUDA \\
			CUDA    & Compute Unified Device Architecture & Nền tảng tính toán song song trên GPU của NVIDIA \\
			CUTLASS & CUDA Templates for Linear Algebra Subroutines & Thư viện mẫu C++ cho đại số tuyến tính trên CUDA \\
			DRAM    & Dynamic Random-Access Memory & Bộ nhớ truy cập ngẫu nhiên động \\
			FFI     & Foreign Function Interface & Giao diện gọi hàm giữa các ngôn ngữ \\
			FP32    & 32-bit Floating Point & Số thực dấu phẩy động 32 bit \\
			GEMM    & General Matrix Multiplication & Phép nhân ma trận tổng quát \\
			GPU     & Graphics Processing Unit & Bộ xử lý đồ họa \\
			IR      & Intermediate Representation & Biểu diễn trung gian \\
			JIT     & Just-In-Time (compilation) & Biên dịch tức thời khi chạy \\
			LLM     & Large Language Model & Mô hình ngôn ngữ lớn \\
			LLVM    & LLVM Compiler Infrastructure (ban đầu là ``Low Level Virtual Machine'') & Hạ tầng trình biên dịch LLVM \\
			MatMul  & Matrix Multiplication & Phép nhân ma trận \\
			REPL    & Read--Eval--Print Loop & Vòng lặp đọc -- thực thi -- in kết quả \\
			\bottomrule
		\end{tabular}
	\end{center}
	
	
	\newpage
	
	% ================= NỘI DUNG CHÍNH =================
	\section{Tổng quan dự án}
	\label{sec:tongquan}
	
	\subsection{Lý do chọn dự án}
	
        Sự phát triển của AI trong thời gian gần đây đã dẫn đến sự bùng nổ của nhiều thuật toán máy học và mô hình học sâu. Do vậy, nhu 
        cầu cho năng lực tính toán để học trên những bộ dữ liệu khổng lồ
        cũng từ đó tăng.
		
		Hiện nay, một trong những phép toán được sử dụng rộng rãi cho thuật toán máy học là GEMM. Trong quá trình học, mỗi phép nhân ma trận sẽ tạo ra một ma trận khác và ma trận mới sẽ được nối với những phép nhân ma trận khác. Để xử lý nhiều thao tác, tác vụ và quá trình khác nhau, những mô hình máy học cần thực hiện hàng triệu phép tính ma trận. Do những phép tính ma trận này so với các phép tính thông thường khác có tính đặc thù về thời gian rất cao, vì vậy những phép tính tiếp theo dần bị dồn lên bộ nhớ đệm và các bước trung gian cho đến khi những chường trình quản lý bộ nhớ hoặc quy tắc của máy tính giải phóng bộ nhớ.
        
        Việc lặp lại quá trình này đối với mô hình máy học trở thành một vấn đề lớn trong việc hoạt động dưới một khối lượng công việc khổng lồ, để lại hậu quá lớn cho hiệu suất và độ trễ của mô hình trên thiết bị và tài nguyên phần cứng nó đang sử dụng.

        Bởi vậy, phần mềm đồng thời góp phần tối ưu hóa nhân của GEMM và tăng tuổi thọ của phần cứng được sử dụng, đảm bảo bộ đệm, dữ liệu đệm và con trỏ được giải phóng sớm nhất có thể và giảm phí tổn từ kết quả của việc dồn bộ nhớ đệm.

	    Dự án được lấy cảm hứng từ những nghiên cứu về hệ thống bộ nhớ đọc tiêu diệt cho ngôn ngữ Sere~\cite{sere}, ngôn ngữ Clean~\cite{clean} và thư viện tối ưu phép nhân ma trận ở GPU.
	
	\subsection{Câu hỏi nghiên cứu}
	
	Tuy nhiên, cần cân nhắc các chi phí phát sinh vốn có, cũng như khả năng đánh giá quá cao mức chi phí bộ nhớ hoặc tính thực tiễn của việc tối ưu bộ nhớ. Do đó, dự án nghiên cứu của nhóm chúng em đưa ra các giả thuyết sau:
	
	\begin{itemize}
	\item Ngữ nghĩa đọc phá hủy có thể giảm mức sử dụng bộ nhớ đỉnh đối với các khối lượng công việc chứa nhiều ma trận trung gian có vòng đời ngắn.
	\item Đọc phá hủy có thể giảm số lượng hoặc chi phí cấp phát bộ đệm tạm thời khi các bộ đệm đã bị tiêu thụ có thể được tái sử dụng một cách an toàn.
	\item Với các khối lượng công việc mà việc quản lý bộ nhớ chiếm tỷ trọng đáng kể trong tổng thời gian thực thi, việc tái sử dụng bộ đệm có thể cải thiện hiệu năng đầu-cuối.
	\end{itemize}

    Đồng thời qua nghiên cứu này có thể trả lời được câu hỏi chính về việc "Làm thế nào để tối ưu hóa việc tính toán các phép nhân ma trận ở trong bộ nhớ đệm và trung gian để giảm chi phí máy tính?".
	
	\subsection{Mục tiêu nghiên cứu}
	
	% TODO: viết mục tiêu nghiên cứu
	
	\subsection{Tính mới nghiên cứu}
	
	Mục tiêu chính của chúng en là liên kết trực tiếp thư viện runtime cuBLAS đã được biên dịch sẵn cho một ứng dụng cụ thể, được đánh giá bằng một phương pháp cụ thể, theo một mô hình ràng buộc tương tự hệ thống kiểu tuyến tính nghiêm ngặt. Một số ngôn ngữ khác cũng đã sử dụng các ngữ nghĩa để theo dõi việc tiêu thụ tương tự: hệ thống sở hữu (ownership) của Rust~\cite{rust} áp đặt nguyên tắc chỉ sử dụng một lần thông qua hệ thống kiểu affine (biến thể nới lỏng của hệ thống kiểu tuyến tính, chỉ sử dụng biến số một lần duy nhất), và Futhark~\cite{futhark} - một ngôn ngữ mảng thuần hàm sử dụng kiểu duy nhất (uniqueness types) để hỗ trợ cập nhật tại chỗ một cách an toàn.
    
    Tuy nhiên, một hạn chế của Futhark trong quá trình vận hành là việc hệ thống tự sinh các kernel (nhân) GPU từ mã nguồn thay vì tích hợp với các thư viện hiệu năng cao bên ngoài được duy trì độc lập như cuBLAS hay CUTLASS, bởi vậy nó được xem là một trong những trở ngại thực tiễn của ngôn ngữ này. Ngược lại, MemoLang tách biệt quá trình phân tích tiêu thụ tại thời điểm biên dịch khỏi quá trình thực thi kernel.
	
	\section{Cơ sở lý thuyết}
	
	% TODO: Trình bày động lực của ngữ nghĩa đọc phá hủy, giải thích vì sao nó
	% giúp sử dụng bộ nhớ hiệu quả cho các khối lượng công việc ma trận, và bổ sung
	% kiến thức nền về sinh LLVM IR hoặc mô hình lập trình của cuBLAS.
	
	Phép nhân ma trận có thể được biểu diễn như sau:
	
	\begin{equation}
		C_{i,j} = \sum_{k=0}^{N-1} A_{i,k}B_{k,j}
	\end{equation}
	
	Mỗi phần tử trong ma trận kết quả $C$ được tạo ra bằng cách nhân một hàng của $A$ với một cột của $B$ rồi cộng các tích lại. Các cài đặt hiệu quả phải quản lý việc di chuyển dữ liệu, tái sử dụng bộ nhớ và, trong cơ chế đọc phá hủy, vòng đời của các bộ đệm toán hạng sau khi chúng đã bị tiêu thụ.
	
	\begin{equation}
	T_1 = A B,\qquad
	T_2 = T_1 C,\qquad
	T_3 = T_2 D,
	\end{equation}
	
	Thực tế, nhiều bộ đệm trung gian cùng tồn tại tại một thời điểm. Vùng nhớ của các bộ đệm này nên được xóa và tái sử dụng ngay sau khi các phép toán tương ứng hoàn tất hoặc khi vòng đời của đối tượng kết thúc để có được hiệu quả tốt nhất.
	
	
	Đọc phá hủy (destructive read) là thao tác đọc bộ nhớ mà bản thân nó xóa đi giá trị được đọc: dữ liệu bị hủy ngay tại thời điểm được truy cập. Nếu cần giữ lại giá trị, một thao tác ghi ngược (write-back) phải khôi phục nó ngay lập tức, hoặc một bộ đệm/bộ nhớ đệm bên ngoài phải lưu giữ một bản sao.
	
	Trong phần mềm hiện đại, kỹ thuật này hiếm khi được sử dụng trực tiếp do độ phức tạp khi cài đặt. Thay vào đó, các ngôn ngữ lập trình hiện đại sử dụng cơ chế sở hữu và tính bất biến để hiện thực hóa cơ chế di chuyển ngữ nghĩa (move semantics), nhằm đạt được mức độ an toàn và khả năng tái sử dụng bộ nhớ tương tự mà không gặp phải những rủi ro của việc đọc phá hủy không đồng bộ. Các nghiên cứu liên quan về hệ thống kiểu tuyến tính cũng tiếp cận tương tự ở cấp độ ngôn ngữ, yêu cầu mỗi biến phải được sử dụng đúng một lần.
	
	Trong lịch sử, đọc phá hủy dần trở nên phổ biến hơn trong thiết kế phần cứng và mạch điện hơn là trong phần mềm. Dybdahl et al.~\cite{DybdahlKGN06} đã nghiên cứu hành vi đọc phá hủy trong DRAM nhúng được sử dụng làm bộ nhớ chính, dựa trên cơ chế sau: ``Dữ liệu trước tiên được đọc từ DRAM vào bộ nhớ đệm. Tại thời điểm này, dữ liệu không được lưu trong DRAM mà chỉ nằm trong bộ nhớ đệm. Dữ liệu được ghi ngược lại DRAM khi dòng bộ nhớ đệm bị thay thế''~\cite{DybdahlKGN06}.
    
    Việc loại bỏ thao tác ghi ngược tức thời theo cách này mang lại lợi ích đo được về điện năng và hiệu năng, bao gồm việc cải thiện thời gian bộ đệm ghi ngược từ 6~ns xuống 3~ns. Mặc dù kết quả này chỉ áp dụng cho một hệ thống bộ nhớ phần cứng cụ thể, nó gợi mở câu hỏi rộng hơn mà dự án này nghiên cứu: vậy một nguyên tắc tương tự, tức tiêu thụ khi đọc, khi được áp dụng ở cấp độ phần mềm/trình thông dịch cho các khối lượng công việc nặng về GEMM, có thể mang lại lợi ích tương tự về áp lực bộ nhớ và chi phí cấp phát hay không.
	
	\section{Mục tiêu dự án}
	
	Mục tiêu của MemoLang là xác định liệu ngữ nghĩa đọc phá hủy ở cấp độ ngôn ngữ có thể mang lại lợi ích đo lường được về quản lý bộ nhớ cho các chương trình GPU xử lý ma trận, đồng thời vẫn giữ được hiệu năng của các kernel nhân ma trận đã được tối ưu hay không.
	
	Dự án sẽ được triển khai và đánh giá thông qua các mục tiêu sau:
	
	\begin{itemize}
		\item Định nghĩa cú pháp cốt lõi tương tự Python và ngữ nghĩa đọc phá hủy của MemoLang.
		\item Xây dựng bộ phân tích từ vựng (lexer), bộ phân tích cú pháp (parser), AST và quy trình sinh mã dựa trên LLVM.
		\item Sinh LLVM IR cho các phép toán số học và phép toán ma trận.
		\item Tích hợp phép nhân ma trận GPU đã được tối ưu thông qua CUTLASS và/hoặc cuBLAS.
		\item Hiện thực cơ chế theo dõi quyền sở hữu và việc tiêu thụ một cách tường minh cho các toán hạng ma trận.
		\item Xác định khi nào các bộ đệm ma trận đã bị tiêu thụ có thể được tái sử dụng an toàn.
		\item Kiểm chứng kết quả sinh ra so với một tham chiếu số học độc lập.
		\item Đo mức sử dụng bộ nhớ GPU đỉnh và hành vi cấp phát bộ đệm tạm thời.
		\item Đo cả thời gian thực thi kernel GPU và tổng thời gian chạy đầu-cuối của chương trình.
		\item So sánh việc thực thi có đọc phá hủy với một cài đặt tương đương không phá hủy.
		\item So sánh MemoLang với các thư viện GPU đã tối ưu làm mốc tham chiếu khi phù hợp.
		\item Lưu trữ nhật ký đo hiệu năng, kết quả kiểm tra tính đúng đắn, đầu ra của trình biên dịch và các sản phẩm build để đảm bảo khả năng tái lập kết quả thực nghiệm.
	\end{itemize}
	
	Quá trình đánh giá sẽ không mặc định rằng đọc phá hủy mang lại cải thiện hiệu năng. Việc giảm mức sử dụng bộ nhớ đỉnh mà không giảm thời gian chạy vẫn được xem là một kết quả có ý nghĩa; còn nếu quan sát thấy thời gian chạy được cải thiện, kết quả đó sẽ được phân tích để xác định liệu nó bắt nguồn từ việc giảm chi phí quản lý bộ nhớ hay từ thay đổi trong kernel nhân ma trận bên dưới.
	
	\subsection{Ngữ nghĩa đọc phá hủy}
	
	Một thao tác đọc phá hủy coi một giá trị là đã bị tiêu thụ khi một phép toán có quyền sử dụng độc quyền giá trị đó. Sau khi bị tiêu thụ, giá trị ban đầu không còn khả dụng cho các phép toán tiếp theo, trừ khi lập trình viên tạo bản sao một cách tường minh hoặc bảo toàn nó bằng cách khác.
	
	Đối với các phép toán ma trận, điều này tạo cơ hội tái sử dụng vùng nhớ. Ví dụ,
	
	\begin{lstlisting}
		C = A @ B
		D = C @ W
	\end{lstlisting}
	
	có thể cho phép vùng nhớ của $A$ hoặc $C$ được tái sử dụng khi trình biên dịch xác định rằng giá trị tương ứng không còn được sử dụng ở đâu nữa.
	
	Mục đích của cơ chế này không phải là thay đổi phép tính toán học của phép nhân ma trận. Thay vào đó, nó thay đổi vòng đời của các bộ đệm liên quan đến phép tính. Nếu một bộ đệm có thể được tái sử dụng mà không cần sao chép nội dung, cài đặt có thể giảm mức tiêu thụ bộ nhớ đỉnh và số lần cấp phát.
	
	Tối ưu này phụ thuộc vào đặc thù của khối lượng công việc. Một ma trận còn được sử dụng ở nhiều nơi về sau thì không thể bị hủy an toàn sau lần sử dụng đầu tiên. Trong trường hợp đó, cài đặt có thể cần sao chép, mượn (borrow) hoặc dùng một cơ chế khác để bảo toàn giá trị. Do đó, ngữ nghĩa đọc phá hủy có thể tạo ra sự đánh đổi giữa hiệu quả sử dụng bộ nhớ và tính linh hoạt khi tái sử dụng giá trị.
	
	Vì vậy, ngôn ngữ coi đọc phá hủy là một cơ chế sở hữu tường minh, thay vì mặc định rằng việc tiêu thụ mọi toán hạng chắc chắn sẽ cải thiện hiệu năng.
	
	\section{Kiến trúc trình biên dịch và môi trường thực thi}
	
	% TODO: Mô tả cú pháp và ngữ nghĩa của MemoLang. Gợi ý các mục con:
	% - Cú pháp bề mặt và các quy ước giống Python
	% - Hệ thống kiểu (nếu có) cho ma trận/tensor
	% - Ngữ nghĩa đọc phá hủy: một lần đọc tiêu thụ giá trị như thế nào, điều gì
	%   xảy ra khi tái sử dụng, mô hình lỗi/quyền sở hữu
	% - Mô hình lập trình CUDA đã ảnh hưởng đến thiết kế như thế nào
	
	\subsection{Thiết kế cú pháp}
	
	Chúng tôi lựa chọn quy ước cú pháp tương tự Python cho ngôn ngữ lập trình:
	\begin{itemize}
		\item Chú thích (comment) bằng dấu thăng (\#).
		\item Các kiểu nguyên thủy: số nguyên (Integer), số thực (Float), luận lý (Boolean) và chuỗi (String).
		\item Khai báo biến.
		\item Định nghĩa hàm bằng từ khóa `def'.
		\item Từ khóa extern cho liên kết runtime/FFI (liên ngôn ngữ).
		\item Độ ưu tiên toán tử.
		\item Cấu trúc điều khiển luồng.
	\end{itemize}
	Tuy nhiên, vì tính thực tiễn khi cài đặt và các quyết định thiết kế khác, chúng tôi đã điều chỉnh một số điểm so với cú pháp Python gốc:
	\begin{itemize}
		\item Độ ưu tiên đặc biệt cho clone() dùng để sao chép vùng nhớ.
		\item Sử dụng dấu ngoặc nhọn thay cho thụt lề để xác định phạm vi của hàm.
	\end{itemize}
	
	\subsection{Quy trình xử lý}
	% TODO: Mô tả quy trình: mã nguồn -> lexer -> parser -> ExprAST ->
	% LLVM IR -> biên dịch JIT/AOT -> gọi cuBLAS cho các kernel nhân ma trận.
	\begin{figure}[H]
	\centering
	\begin{tikzpicture}[
	node distance=5cm,
	box/.style={rectangle, draw, rounded corners, minimum width=6cm, text width=5.6cm, minimum height=3cm, align=center, font=\normalsize},
	special/.style={box, fill=gray!20},
	arr/.style={-{Latex[length=1.5mm]}}
	]
	\node[box]    (src) at (0,0)  {Mã nguồn \\ \small Văn bản MemoLang thô};
	\node[box]   (lex) at (10,0)  {Bộ phân tích từ vựng \\ \small Tạo chuỗi token};
	\node[box]   (par) at (10,-5) {Bộ phân tích cú pháp \\ \small Xây dựng AST theo BinopPrecedence};
	\node[box]   (ast) at (0,-5) {AST \\ \small Cây nút / biểu thức / câu lệnh};
	\node[special] (con) at (0,-10)  {Phân tích tiêu thụ [Đọc phá hủy] \\ \small Đánh dấu consume\_a / consume\_b};
	\node[box]   (ir) at (10,-10)   {LLVM IR \\ \small Sinh lời gọi kèm cờ tiêu thụ};
	\node[box]    (jit) at (10,-15)  {JIT / AOT \\ \small Biên dịch sang mã máy};
	\node[box]   (link) at (0,-15) {Liên kết \\ \small Liên kết runtime của MemoLang};
	\node[special] (gpu) at (0,-20)  {Kernel GPU [Đọc phá hủy] \\ \small Giải phóng vùng nhớ đã tiêu thụ};
	\draw[arr] (src) -- (lex);
	\draw[arr] (lex) -- (par);
	\draw[arr] (par) -- (ast);
	\draw[arr] (ast) -- (con);
	\draw[arr] (con) -- (ir);
	\draw[arr] (ir) -- (jit);
	\draw[arr] (jit) -- (link);
	\draw[arr] (link) -- (gpu);
	\end{tikzpicture}
	\caption{Quy trình xử lý bộ nhớ, bao gồm toàn bộ quá trình từ văn bản mã nguồn đến thực thi trên GPU. Các khối tô xám biểu thị những bước liên quan đến ngữ nghĩa đọc phá hủy.}
	\label{fig:pipeline}
	\end{figure}
	
	\section{Toán học và thuật toán}
	
	\subsection{Cơ sở toán học của phép nhân ma trận FP32}
	
	% TODO: Trình bày chi tiết về tích lũy FP32, chia khối (tiling) và các vấn đề
	% về độ chính xác của kernel GEMM.
	
	\subsection{Cơ chế giải phóng khi tiêu thụ}
	
	% TODO: Mô tả hình thức hoặc phi hình thức về đọc phá hủy: khi một giá trị ma
	% trận đã được đọc (ví dụ bị tiêu thụ bởi @), nó không thể được dùng lại trừ khi
	% được sao chép tường minh. Giải thích cách điều này giảm bộ nhớ đỉnh, và cách các
	% cờ consume_a/consume_b được tính toán và truyền cho môi trường thực thi.
	
	\section{Tích hợp CUTLASS}
	
	% TODO: Mô tả quy trình: mã nguồn -> lexer -> parser -> ExprAST ->
	% LLVM IR -> biên dịch JIT/AOT -> gọi CUTLASS cho các kernel nhân ma trận.
	
	\begin{table}[h]
		\centering
		\caption{Các tệp mã nguồn chính}
		\label{tab:modules}
		\begin{tabular}{ll}
			\toprule
			Tệp & Chức năng \\
			\midrule
			\texttt{Lexer.h/.cpp} & Tách mã nguồn MemoLang thành token \\
			\texttt{Parser.h/.cpp} & Xây dựng AST \\
			\texttt{ExprAST.h} & Định nghĩa các nút AST \\
			\texttt{CodeGen.cpp} & Sinh LLVM IR \\
			\texttt{CuBLASRuntime.cpp} & Liên kết cuBLAS cho phép nhân ma trận \\
			\texttt{main.cpp} & Chương trình điều khiển / REPL \\
			\bottomrule
		\end{tabular}
	\end{table}
	
	% TODO: Mô tả cách biểu thức nhân ma trận được hạ xuống LLVM IR, và khi nào/như
	% thế nào quyền điều khiển được chuyển cho CUTLASS (ví dụ qua một lời gọi hàm bên
	% ngoài được chèn khi sinh mã). Trình bày việc bàn giao bộ đệm đọc phá hủy cho CUTLASS.
	
	\section{Kiểm chứng}
	
	% TODO: Mô tả cách kiểm tra tính đúng đắn (ví dụ so sánh kết quả nhân ma trận
	% của MemoLang với kết quả tham chiếu từ NumPy hoặc cuBLAS cho nhiều kích thước
	% ma trận). Thay bằng bảng kết quả khi có.
	
	\begin{table}[h]
		\centering
		\caption{Kết quả kiểm tra tính đúng đắn (mẫu)}
		\label{tab:verification}
		\begin{tabular}{ll}
			\toprule
			Trường hợp kiểm tra & Kết quả \\
			\midrule
			Nhân ma trận $N \times N$ so với tham chiếu & Chưa có \\
			Phát hiện lỗi tái sử dụng sau đọc phá hủy & Chưa có \\
			Trường hợp biên (rỗng, 1x1, không vuông) & Chưa có \\
			\bottomrule
		\end{tabular}
	\end{table}
	
	\section{Đánh giá hiệu năng}
	
	Việc đánh giá hiệu năng được thiết kế để phân biệt cải thiện về quản lý bộ nhớ với cải thiện của bản thân phép nhân ma trận.
	
	Ba cài đặt chính sẽ được so sánh:
	
	\begin{enumerate}
		\item Một cài đặt sử dụng thư viện GPU đã tối ưu, làm mốc tham chiếu cho hiệu năng kernel.
		\item MemoLang với ngữ nghĩa giá trị thông thường (không phá hủy).
		\item MemoLang với ngữ nghĩa đọc phá hủy.
	\end{enumerate}
	
	Khi phù hợp, một cài đặt C++ hoặc CUDA đơn giản cũng sẽ được đưa vào làm mốc cơ sở.
	
	Phép so sánh quan trọng nhất là giữa hai cấu hình MemoLang. Cả hai cấu hình thực hiện cùng các phép toán và sử dụng cùng một cài đặt nhân ma trận GPU bên dưới. Phép so sánh này tách biệt tác động của ngữ nghĩa đọc phá hủy hiệu quả hơn so với việc so sánh với một cài đặt không liên quan.
	
	Bộ đánh giá sẽ bao gồm cả các phép nhân ma trận riêng lẻ lẫn các chuỗi phép toán ma trận. Kích thước ma trận gồm nhỏ, trung bình, lớn và không vuông. Điều này cho phép xác định liệu tác động của đọc phá hủy có phụ thuộc vào quy mô khối lượng công việc hay không.
	
	Các số liệu sau sẽ được thu thập:
	
	\begin{itemize}
		\item Thời gian thực thi kernel GPU.
		\item Thời gian thực thi đầu-cuối.
		\item Mức tiêu thụ bộ nhớ GPU đỉnh.
		\item Số lần cấp phát bộ đệm tạm thời.
		\item Số lần tái sử dụng bộ đệm.
		\item Thời gian truyền dữ liệu từ máy chủ sang thiết bị và từ thiết bị về máy chủ (nếu có).
		\item Chi phí của trình biên dịch và môi trường thực thi.
		\item Sai số số học so với cài đặt tham chiếu.
	\end{itemize}
	
	Các phép đo trên GPU sẽ được lặp lại nhiều lần với số vòng khởi động (warm-up) phù hợp. Chi phí khởi tạo CUDA và chi phí biên dịch sẽ được tách riêng khỏi các phép đo ở trạng thái ổn định khi có thể.
	
	Kết quả đánh giá sẽ báo cáo cả giá trị tuyệt đối lẫn mức thay đổi tương đối. Không có cải thiện hiệu năng nào được quy cho ngữ nghĩa đọc phá hủy trừ khi các số liệu đo được ủng hộ cách diễn giải đó.
	


	\subsection{Mối quan hệ kỳ vọng giữa ngôn ngữ và hiệu năng nhân ma trận}
	
	Có thể kỳ vọng một ngôn ngữ với ngữ nghĩa đọc phá hủy sẽ mang lại lợi ích về quản lý bộ nhớ. Trong khi đó, các phép toán số học như phép nhân ma trận có thể được giao cho kernel CUTLASS/cuBLAS đã tối ưu nhằm khắc phục hạn chế của Futhark~\cite{futhark}. Việc cài đặt được kỳ vọng không ảnh hưởng đến hiệu năng tính toán so với cài đặt trên các ngôn ngữ khác. Tuy nhiên, hiệu năng đầu-cuối được kỳ vọng sẽ cải thiện nhờ các tối ưu của ngôn ngữ, bao gồm:
	\begin{itemize} 
		\item Giảm cấp phát tạm thời
		\item Giảm vòng đời bộ đệm
		\item Giảm áp lực lên bộ nhớ GPU
	\end{itemize}
	Về lý thuyết, hiệu quả của cài đặt này tăng theo quy mô khối lượng công việc, vì các tensor tức thời hoặc các phép toán ma trận nhỏ không chịu tác động đáng kể từ việc quản lý bộ nhớ so với các chuỗi công việc lớn.
	
	Thời gian thực thi kỳ vọng có thể được phân tách gần đúng như sau:
	
	\begin{equation}
		T_{\mathrm{total}} =
		T_{\mathrm{language}} +
		T_{\mathrm{memory}} +
		T_{\mathrm{launch}} +
		T_{\mathrm{kernel}}.
	\end{equation}
	
	\noindent trong đó $T_{\mathrm{language}}$ là chi phí của ngôn ngữ, $T_{\mathrm{memory}}$ là chi phí quản lý bộ nhớ, $T_{\mathrm{launch}}$ là chi phí khởi chạy kernel và $T_{\mathrm{kernel}}$ là thời gian thực thi kernel.
	
	Ngữ nghĩa đọc phá hủy chủ yếu nhắm vào $T_{\mathrm{memory}}$. Chúng không được kỳ vọng sẽ giảm đáng kể $T_{\mathrm{kernel}}$ khi cùng một kernel GPU đã tối ưu được sử dụng.
	
	Với một phép nhân ma trận lớn, $T_{\mathrm{kernel}}$ có thể chiếm phần lớn tổng thời gian thực thi, dẫn đến cải thiện đầu-cuối khó quan sát được. Tuy nhiên, với các khối lượng công việc chứa nhiều tensor trung gian hoặc các phép toán ma trận nhỏ, $T_{\mathrm{memory}}$ và các chi phí khác của môi trường thực thi có thể trở nên đáng kể hơn.
	
	Do đó, một kết quả thành công không đòi hỏi MemoLang phải làm cho từng phép nhân ma trận riêng lẻ nhanh hơn. Việc giảm mức sử dụng bộ nhớ đỉnh, số lần cấp phát hoặc chi phí đầu-cuối vẫn có thể là một kết quả có ý nghĩa, ngay cả khi thời gian thực thi kernel gần như không đổi.
	
	
	
		
	\section{Minh chứng và khả năng tái lập}
	\label{sec:evidence}
	
	% TODO: Liệt kê nơi lưu nhật ký, kết quả đo hiệu năng và các sản phẩm build
	% để mọi khẳng định trong báo cáo đều có thể truy vết tới dữ liệu đã lưu.
	
	
	\section{Kết luận}
	
	% TODO: Tóm tắt những gì đã xây dựng và kiểm chứng, những gì còn là hướng phát
	% triển, gắn các khẳng định với minh chứng đã lưu ở Phần~\ref{sec:evidence}.
	
	Các lợi ích thực tiễn tiềm năng bao gồm:
	
	\begin{itemize}
		\item Cải thiện hiệu năng các phép nhân ma trận, hỗ trợ huấn luyện và suy luận nhanh hơn cho các mô hình học máy.
		\item Trình biên dịch tái sử dụng hiệu quả hơn các bộ đệm tensor trung gian.
		\item Giảm chi phí cấp phát bộ đệm và bộ nhớ đệm khi chạy.
	\end{itemize}
	
	
	\subsection{Tiến độ}
	
	% TODO: mô tả tiến độ hiện tại của dự án
	
	\addcontentsline{toc}{section}{\refname}
	\begin{thebibliography}{9}
		\bibitem{DybdahlKGN06}
		Haakon Dybdahl, Per Gunnar Kjeldsberg, Marius Grann{\ae}s, and Lasse Natvig,
		``Destructive-read in embedded DRAM, impact on power consumption,'' \url{https://www.researchgate.net/publication/220222072_Destructive-read_in_embedded_DRAM_impact_on_power_consumption}
		
		\bibitem{cublas}
		NVIDIA, ``cuBLAS Library.'' \url{https://developer.nvidia.com/cublas}
		
		\bibitem{llvm}
		LLVM Project, ``The LLVM Compiler Infrastructure.'' \url{https://llvm.org}
		
		\bibitem{clean}
		Clean Language Team, ``The Clean Programming Language.'' \url{https://clean.cs.ru.nl}
		
		\bibitem{sere}
		Sere Language Project, ``Sere Programming Language.'' \url{https://clean-lang.org}
		
		\bibitem{rust}
		N.~D.~Matsakis and F.~S.~Klock, ``The Rust language,'' in \emph{Proc. ACM SIGAda Ada Letters}, 2014. \url{https://www.rust-lang.org}
		
		\bibitem{futhark}
		T.~Henriksen, N.~G.~W.~Serup, M.~Elsman, F.~Henglein, and C.~E.~Oancea, ``Futhark: purely functional GPU-programming with nested parallelism and in-place array updates,'' in \emph{Proc. 38th ACM SIGPLAN Conference on Programming Language Design and Implementation (PLDI)}, 2017. \url{https://futhark-lang.org}
		
		\bibitem{memolang}
		Independent Interpreter Project, ``MemoLang: mã nguồn và minh chứng đo hiệu năng,'' kho mã nguồn dự án, tháng 9 năm 2026.
		
	\end{thebibliography}
	
\end{document}